\section{Animación procedural}
La dinámica principal del sistema reside en la animación del flujo de datos entre capas consecutivas, 
representada mediante la ejecución procedural de la operación de convolución.
\medskip

La animación de la convolución se implementa en el grupo de nodos \textit{Volume Connection}, modelando 
la iteración espacial del núcleo sobre las celdas de una cuadrícula bidimensional. Este nodo recibe como 
parámetros de entrada las referencias a dos objetos volumétricos (instancias del grupo de nodos 
\textit{Volume Viewer}), los cuales fungen como capas de entrada y salida respectivamente. Para cada una de estas 
capas, el árbol de nodos ejecuta de forma desacoplada la siguiente secuencia de pasos:
\begin{enumerate}
    \item \textbf{Cálculo de parámetros espaciales iniciales:} a partir de la resolución del mapa de activación y sus 
    dimensiones físicas en el espacio 3D, determina el tamaño geométrico del \textit{kernel}, las dimensiones 
    de cada celda sobre el plano de la rebanada y la coordenada de origen de la iteración (esquina superior izquierda).
    \item \textbf{Determinación de la cuadrícula de muestreo:} computa el número total de celdas sobre las 
    cuales itera el \textit{kernel}. En el caso particular de la capa de entrada, el cómputo integra además 
    los valores de paso (\textit{stride}) y el modo de redondeo (\textit{Ceiling Mode}) para definir la 
    grilla discreta de convolución.
    \item \textbf{Muestreo temporal del índice activo:} mediante el parámetro escalar $T \in [0, 1]$, que 
    gobierna el progreso temporal de la animación, calcula el índice secuencial de la celda activa en el 
    instante actual. Para $T = 0$, el índice apunta a la celda inicial (esquina superior izquierda), mientras 
    que para $T = 1$ corresponde a la celda final (esquina inferior derecha).
    \item \textbf{Conversión a coordenadas bidimensionales de grilla:} con base en el índice lineal obtenido y 
    la resolución espacial de la capa, calcula las coordenadas discretas $(X, Y)$ de la celda activa dentro de 
    la cuadrícula de iteración.
    \item \textbf{Transformación y traslación geométrica en el espacio 3D:} utilizando las dimensiones físicas 
    de las celdas obtenidas en el primer paso, evalúa el vector de desfase respecto a la posición de origen y 
    traslada la geometría del \textit{kernel} a su ubicación tridimensional definitiva en la escena.
\end{enumerate}
\medskip

Este procedimiento algorítmico constituye el núcleo de la animación procedural. Dado que el factor $T$ 
se encuentra expuesto como un atributo en la interfaz del grupo de nodos, la convolución puede interpolarse, 
sincronizarse o controlarse directamente mediante fotogramas clave (\textit{keyframes}) en la línea de tiempo de Blender.